\documentclass[10pt,a4paper]{amsart}
\usepackage[margin=19mm]{geometry}
\usepackage{amsmath,amssymb,mathtools}
\usepackage[scr=rsfs,cal=euler]{mathalfa}
\usepackage{xfrac}
\usepackage[unicode,hidelinks,bookmarks=false]{hyperref}
\newcommand{\ie}{i.e.\ }
\newcommand{\cf}{cf.\ }
\newcommand{\resp}{respectively\ }
\newcommand{\Subsection}{\subsection}
\providecommand{\nobreakdash}{\nobreak\hbox{-}}
\setlength{\emergencystretch}{2em}
\multlinegap0pt
\usepackage{amssymb}
\usepackage{mathtools}
\usepackage{tikz}
\usepackage[shortlabels]{enumitem}
\usepackage{xcolor}
\usepackage{dsfont}
\usepackage[normalem]{ulem}
\newtheorem{theorem}{Theorem}
\newtheorem{proposition}{Proposition}
\newtheorem{lemma}{Lemma}
\def\a{\alpha}
\def\d{\delta}
\title{Some remarks on Folkman graphs for triangles}
\author{Eion Mulrenin, Steven Van Overberghe}
\date{Source: arXiv:2506.14942v4 (CC BY 4.0)}
\begin{document}
\maketitle

\noindent\emph{Source and licence.} Some remarks on Folkman graphs for triangles. Version arXiv:2506.14942v4 (CC BY 4.0), distributed by arXiv under the Creative Commons Attribution 4.0 International licence, http://creativecommons.org/licenses/by/4.0/, as recorded in the arXiv licence metadata for this exact version and shown on the abstract page https://arxiv.org/abs/2506.14942v4. Full licence notice: \texttt{licenses/arxiv-2506.14942-CC-BY-4.0.txt}.\par\smallskip

\noindent\textit{Editorial scope.} Counted unit: the article's theorem main-2 (source0.tex lines 260--264). The article's complete original proof of it is delivered with it (source0.tex lines 784--809, one \texttt{proof} environment of 26 lines, which the article places away from the statement). No mathematics is written, repaired or re-derived: the statement and the proof are the article's own lines, in the article's own notation, and no retained line is substituted or re-flowed.  Retained uncounted as the source-local context the proof needs: lines 634--642; lines 669--673; lines 710--716.  Nothing was omitted from any retained span, so the package carries no omission record.  No retained line contains a citation, so the wrapper carries no bibliography and no bibliography entry.  No retained line contains a figure environment, an image-inclusion command or drawing code, so no image is delivered and the package declares no asset dependency.  Editorial formatting only, in addition: an AMS \texttt{amsart} preamble that restates the article's own declarations and macros used by the retained lines, namely the environments \texttt{theorem}, \texttt{proposition}, \texttt{lemma} on separate counters, together with the standard packages those lines need.  The retained environments print in this document as Theorem 1, Proposition 1, Lemma 1 in order of appearance, whereas the article prints the same items with the numbering of its own full document.  The complete native source of the article as distributed in the arXiv e-print for this exact version is bundled unchanged as \texttt{sources/180047/source0.tex}.
\begin{proposition}
\label{prp: goodman}
    Let $G = (V,E)$ be a graph, and let $\Delta: E \to \{R,B\}$ be a two-coloring of its edges.
    For a vertex $v \in V$, let $R_{\Delta}(v)$ ($B_\Delta(v)$) denote the number of edges in $N(v)$ in which both endpoints form a red (blue) edge with $v$.
    Then the number of monochromatic triangles in $G$ under $\Delta$ is given by
    \begin{equation*}
        \frac{1}{2} \sum_{v \in V} {\Big (} R_\Delta(v) + B_\Delta(v) - \frac{1}{3}|E(G[N(v)])| {\Big )}.
    \end{equation*}
\end{proposition}

\begin{lemma}
\label{lem: blowup-coloring}
    Let $F[t]$ be a $t$-blowup of $F$ with $nt$ vertices and $mt^2$ edges.
    Under any coloring $\chi: V(F[t]) \to \{R,B\}$ of its vertex set, there are at least $(1 - \a)mt^2$ monochromatic edges.
\end{lemma}

\begin{lemma}
\label{lem: blowup}
    For every $\d > 0$, the following event holds with high probability, i.e., with probability approaching 1 as $q \to \infty$.
    For every $v \in V(H_q)$ and every spanning clique $C$ in $G_v$,
    %of size $q+1$, 
    each of the $n$ blown-up vertices in $C$ will have\footnote{The notation $(1 \pm \d) \frac{2m(q+1)}{n^3}$ means that the number lies between $(1 - \d) \frac{2m(q+1)}{n^3}$ and $(1 + \d) \frac{2m(q+1)}{n^3}$.} $(1 \pm \d) \frac{2m(q+1)}{n^3}$ vertices in $N^*(v)$.
\end{lemma}

\begin{theorem}
\label{thm: main-2}
    There exists a prime power $q_0$ such that the following holds.
    For all prime powers $q \geq q_0$, one may delete a subset of edges of $H_q$ to make it $K_4$-free, but preserve the property that any two-coloring of the remaining edges yields a monochromatic $K_3$. 
\end{theorem}

\begin{proof}[Proof of Theorem~\ref{thm: main-2}]
Fix $\d > 0$ sufficiently small, fix an instance $H_q^*$ in which Lemma~\ref{lem: blowup} holds, and fix a vertex $v \in V(H_q^*)$.
Note that since $F$ is triangle-free, any blowup of $F$ is also triangle-free, and thus as explained above, $H_q^*$ is $K_4$-free.
Hence, we need to check that the Ramsey property still holds in $H_q^*$, for which we use Proposition~\ref{prp: goodman}.

To that end, fix 
%an edge-coloring $\Delta$ of $H_q^*$ and 
a vertex $v \in V(H_q^*)$, and recall that $H_q^*[N^*(v)] = G_v^*$.
For each clique $C$ in $G_v$, in $G_v^*$ there now exists $m \cdot \left((1\pm\d)\frac{2m(q+1)}{n^3}\right)^2$ edges, and so in total, the number of edges in $G_v^*$ is at most $m \cdot (q^3-q) \cdot \left((1+\d)\frac{2m(q+1)}{n^3}\right)^2$.
However, by Lemma~\ref{lem: blowup-coloring}, it follows that for any choice of an edge-coloring $\Delta$ of $H_q^*$, each $C$ will contribute at least 
$(1 - \a) m \cdot \left((1-\d)\frac{2m(q+1)}{n^3}\right)^2$
to the sum $\sum_{v \in V(H_q^*)} R_\Delta(v) + B_\Delta(v)$, 
%under any coloring $\Delta$, 
and so
%, as there are $q^3-q$ such cliques $C$ in $G_v$, for any coloring $\Delta$, 
we will have
\begin{equation*}
    \sum_{v \in V(H_q^*)} R_\Delta(v) + B_\Delta(v) \geq |V(H_q^*)|(1-\a) m \cdot (q^3-q) \cdot \left((1-\d)\frac{2m(q+1)}{n^3}\right)^2.
\end{equation*}
%and in fact this bound is independent of $\Delta$.
Therefore, by Proposition~\ref{prp: goodman} the number of monochromatic triangles in $H_q^*$ under any coloring $\Delta$ is at least
\begin{equation*}
    \frac{1}{2}|V(H_q^*)| \left( \left(1-\a\right) m \cdot (q^3-q) \cdot \left((1-\d)\frac{2m(q+1)}{n^3}\right)^2 - \frac{1}{3} m \cdot (q^3-q) \cdot \left((1+\d)\frac{2m(q+1)}{n^3}\right)^2 \right),
\end{equation*}
which is positive for $\d$ small enough by the hypothesis that $\a < 2/3$.
\end{proof}

\end{document}
